
\subsection{弗洛尔上同调不变性证明的收束}
\label{subsec:5-3-7}
\liftIndex{I204}

现在把本节与前一节的结果结合起来，
完成定理 4.1.4 和 4.1.5 的证明。

\begin{proof}[定理 4.1.4 的证明]
沿用定理 4.1.4 的记号。
由定理 5.3.1 和引理 5.2.19，
得到严格滤过 $A_\infty$ 双模同态
\begin{equation}
\begin{aligned}
{}^{b^{(1)}}\varphi^{b^{(0)}}\colon\
&\bigl(
C_1(L^{(1)},L^{(0)};\Lambda_{0,\mathrm{nov}}),
{}^{b^{(1)}}\mathfrak n^{b^{(0)}}
\bigr)\\
&\longrightarrow
\bigl(
C_2(L^{(1)\prime},L^{(0)\prime};
\Lambda_{0,\mathrm{nov}}),
{}^{\psi_*b^{(1)}}\mathfrak n^{\psi_*b^{(0)}}
\bigr).
\end{aligned}
\tag{5.3.66}\label{eq:5-3-66}
\end{equation}
它诱导上链同态
\begin{equation}
\begin{aligned}
{}^{b^{(1)}}\varphi^{b^{(0)}}\colon\
&\bigl(
C_1(L^{(1)},L^{(0)};\Lambda_{0,\mathrm{nov}}),
\delta_{b^{(1)},b^{(0)}}
\bigr)\\
&\longrightarrow
\bigl(
C_2(L^{(1)\prime},L^{(0)\prime};
\Lambda_{0,\mathrm{nov}}),
\delta_{\psi_*b^{(1)},\psi_*b^{(0)}}
\bigr),
\end{aligned}
\tag{5.3.67}\label{eq:5-3-67}
\end{equation}
从而诱导
\begin{equation}
\begin{aligned}
\bigl({}^{b^{(1)}}\varphi^{b^{(0)}}_{0,0}\bigr)_*
\colon\
&HF\bigl((L^{(1)},b^{(1)}),
(L^{(0)},b^{(0)});\Lambda_{0,\mathrm{nov}}\bigr)\\
&\longrightarrow
HF\bigl((L^{(1)\prime},\psi_*b^{(1)}),
(L^{(0)\prime},\psi_*b^{(0)});
\Lambda_{0,\mathrm{nov}}\bigr).
\end{aligned}
\tag{5.3.68}\label{eq:5-3-68}
\end{equation}
我们将证明 (5.3.68) 是同构。
注意，在以
$C_2(L^{(1)\prime},L^{(0)\prime};
\Lambda_{0,\mathrm{nov}})$
为靶的同态意义下，有
\[
\varphi_*
\equiv
\bigl(\psi|_{L^{(1)}\cap L^{(0)}}\bigr)_*
\colon
C_1(L^{(1)},L^{(0)};\Lambda_{0,\mathrm{nov}})
\pmod{\Lambda^+_{0,\mathrm{nov}}}.
\]

取 $\varphi$ 的同伦逆
\[
\varphi^{-1}\colon
\bigl(C_2(L^{(1)\prime},L^{(0)\prime};
\Lambda_{0,\mathrm{nov}}),\mathfrak n\bigr)
\longrightarrow
\bigl(C_1(L^{(1)},L^{(0)};
\Lambda_{0,\mathrm{nov}}),\mathfrak n\bigr).
\]
一般而言，$\varphi^{-1}$ 并不严格。
\translatornote{原书实体页 337
把下列 $\mathfrak C^{(i)}$
说成 $L^{(i)\prime}$ 一侧的区间模型，
又把 $\Phi$ 的源写成 $C_2$。
但 $\Psi^{(i)}$ 的两个端点及
(5.3.69)、(5.3.71) 的源
都在原侧 $C(L^{(i)})$ 或 $C_1$；
这里按端点型别统一。
见疑误 A-0311、A-0312。}
令 $(\mathfrak D,\mathfrak n)$
为
$[0,1]\times
(C_1(L^{(1)},L^{(0)};
\Lambda_{0,\mathrm{nov}}),\mathfrak n)$
的一个模型；
它是
$(\mathfrak C^{(1)},\mathfrak m)$、
$(\mathfrak C^{(0)},\mathfrak m)$
上的双模。
这里
$(\mathfrak C^{(i)},\mathfrak m)$
是
$[0,1]\times
(C(L^{(i)}),\mathfrak m)$
的模型。
于是存在
\[
\Phi\colon
\bigl(
C_1(L^{(1)},L^{(0)};
\Lambda_{0,\mathrm{nov}}),\mathfrak n
\bigr)
\longrightarrow
(\mathfrak D,\mathfrak n),
\]
满足
\begin{equation}
\mathrm{Eval}_{s=0}\circ\Phi
=\mathrm{id},
\qquad
\mathrm{Eval}_{s=1}\circ\Phi
=\varphi^{-1}\circ\varphi.
\tag{5.3.69}\label{eq:5-3-69}
\end{equation}
$\Phi$ 是覆盖
\[
\Psi^{(i)}\colon
(C(L^{(i)}),\mathfrak m)
\longrightarrow
(\mathfrak C^{(i)},\mathfrak m)
\]
的双模同态。
这里 $\Psi^{(i)}$
是从恒等同态到
$(\psi^{-1})_*\circ\psi_*
\colon(C(L^{(i)}),\mathfrak m)
\to(C(L^{(i)}),\mathfrak m)$
的同伦。
记
\[
\psi^{-1}\colon
(C(L^{(i)\prime}),\mathfrak m)
\longrightarrow
(C(L^{(i)}),\mathfrak m)
\]
为 $\psi$ 的一个同伦逆。

因此，由引理 5.2.18 得到
下列严格滤过 $A_\infty$ 双模同态。
\liftIndex{I205}
这里以及下文，
把覆盖严格滤过 $A_\infty$ 代数同态的
滤过 $A_\infty$ 双模同态称为\emph{严格}。
\begin{equation}
\begin{aligned}
{}^{\psi_*b^{(1)}}
(\varphi^{-1})^{\psi_*b^{(0)}}\colon\
&\bigl(
C_2(L^{(1)\prime},L^{(0)\prime};
\Lambda_{0,\mathrm{nov}}),
{}^{\psi_*b^{(1)}}\mathfrak n^{\psi_*b^{(0)}}
\bigr)\\
&\longrightarrow
\bigl(
C_1(L^{(1)},L^{(0)};
\Lambda_{0,\mathrm{nov}}),
{}^{(\psi^{-1})_*\psi_*b^{(1)}}
\mathfrak n^{(\psi^{-1})_*\psi_*b^{(0)}}
\bigr),
\end{aligned}
\tag{5.3.70}\label{eq:5-3-70}
\end{equation}
\begin{equation}
\begin{aligned}
{}^{b^{(1)}}\Phi^{b^{(0)}}\colon\
&\bigl(
C_1(L^{(1)},L^{(0)};
\Lambda_{0,\mathrm{nov}}),
{}^{b^{(1)}}\mathfrak n^{b^{(0)}}
\bigr)\\
&\longrightarrow
\bigl(
\mathfrak D,
{}^{\Psi_*^{(1)}b^{(1)}}
\mathfrak n^{\Psi_*^{(0)}b^{(0)}}
\bigr).
\end{aligned}
\tag{5.3.71}\label{eq:5-3-71}
\end{equation}
由 (5.3.69)，
$\mathrm{Eval}_{s=s_0}$
诱导严格滤过 $A_\infty$ 双模同态
\begin{equation}
\begin{aligned}
{}^{\Psi_*^{(1)}b^{(1)}}
(\mathrm{Eval}_{s=0})^{\Psi_*^{(0)}b^{(0)}}
\colon\
&\bigl(
\mathfrak D,
{}^{\Psi_*^{(1)}b^{(1)}}
\mathfrak n^{\Psi_*^{(0)}b^{(0)}}
\bigr)\\
&\longrightarrow
\bigl(
C_1(L^{(1)},L^{(0)};
\Lambda_{0,\mathrm{nov}}),
{}^{b^{(1)}}\mathfrak n^{b^{(0)}}
\bigr),
\end{aligned}
\tag{5.3.72}\label{eq:5-3-72}
\end{equation}
\begin{equation}
\begin{aligned}
{}^{\Psi_*^{(1)}b^{(1)}}
(\mathrm{Eval}_{s=1})^{\Psi_*^{(0)}b^{(0)}}
\colon\
&\bigl(
\mathfrak D,
{}^{\Psi_*^{(1)}b^{(1)}}
\mathfrak n^{\Psi_*^{(0)}b^{(0)}}
\bigr)\\
&\longrightarrow
\bigl(
C_1(L^{(1)},L^{(0)};
\Lambda_{0,\mathrm{nov}}),
{}^{(\psi^{-1})_*\psi_*b^{(1)}}
\mathfrak n^{(\psi^{-1})_*\psi_*b^{(0)}}
\bigr).
\end{aligned}
\tag{5.3.73}\label{eq:5-3-73}
\end{equation}

此外，
\begin{equation}
\begin{aligned}
\bigl(
{}^{\Psi_*^{(1)}b^{(1)}}
(\mathrm{Eval}_{s=0})^{\Psi_*^{(0)}b^{(0)}}
\bigr)
\circ
\bigl({}^{b^{(1)}}\Phi^{b^{(0)}}\bigr)
&=\mathrm{id},\\
\bigl(
{}^{\Psi_*^{(1)}b^{(1)}}
(\mathrm{Eval}_{s=1})^{\Psi_*^{(0)}b^{(0)}}
\bigr)
\circ
\bigl({}^{b^{(1)}}\Phi^{b^{(0)}}\bigr)
&=
{}^{\psi_*b^{(1)}}
(\varphi^{-1})^{\psi_*b^{(0)}}
\circ
\bigl({}^{b^{(1)}}\varphi^{b^{(0)}}\bigr).
\end{aligned}
\tag{5.3.74}\label{eq:5-3-74}
\end{equation}
由命题 5.2.38 的第 \textup{(2)} 项，
\({}^{\Psi_*^{(1)}b^{(1)}}
(\mathrm{Eval}_{s=0})^{\Psi_*^{(0)}b^{(0)}}\)
诱导同构
\[
H\bigl(
\mathfrak D;
\delta_{\Psi_*^{(1)}b^{(1)},
\Psi_*^{(0)}b^{(0)}}
\bigr)
\cong
HF\bigl(
(L^{(1)},b^{(1)}),
(L^{(0)},b^{(0)});
\Lambda_{0,\mathrm{nov}}
\bigr).
\]
因此，由 (5.3.74) 第一式，
\({}^{b^{(1)}}\Phi^{b^{(0)}}\)
诱导同构
\begin{equation}
\begin{aligned}
HF\bigl(
(L^{(1)},b^{(1)}),
(L^{(0)},b^{(0)});
\Lambda_{0,\mathrm{nov}}
\bigr)
\cong
H\bigl(
\mathfrak D;
\delta_{\Psi_*^{(1)}b^{(1)},
\Psi_*^{(0)}b^{(0)}}
\bigr).
\end{aligned}
\tag{5.3.75}\label{eq:5-3-75}
\end{equation}
同理，由命题 5.2.38 的第 \textup{(2)} 项，
\({}^{\Psi_*^{(1)}b^{(1)}}
(\mathrm{Eval}_{s=1})^{\Psi_*^{(0)}b^{(0)}}\)
也诱导同构
\begin{equation}
\begin{aligned}
H\bigl(
\mathfrak D;
\delta_{\Psi_*^{(1)}b^{(1)},
\Psi_*^{(0)}b^{(0)}}
\bigr)
\cong
HF\bigl(
&(L^{(1)},
(\psi^{-1})_*\psi_*b^{(1)}),\\
&(L^{(0)},
(\psi^{-1})_*\psi_*b^{(0)});
\Lambda_{0,\mathrm{nov}}
\bigr).
\end{aligned}
\tag{5.3.76}\label{eq:5-3-76}
\end{equation}
由 (5.3.75)、(5.3.76)
及 (5.3.74) 第二式，
可知 (5.3.68) 是单射。
类似地利用从恒等同态到
$\varphi\circ\varphi^{-1}$
的同伦，可证明 (5.3.68) 是满射。
再以类似论证结合定理 5.3.11
和引理 5.2.19，
可知 (5.3.68) 的同构
只依赖 $\psi_s$ 的同痕类。
\end{proof}

\begin{proof}[定理 4.1.5 的证明]
与上面的证明相同，
利用定理 5.3.14 和引理 5.2.19
即可证明定理 4.1.5。
\end{proof}

下面讨论定理 4.1.4 和 4.1.5
的“弱”及／或“形变”版本。

\begin{theorem}{4.1.4a}\label{thm:4-1-4a}
设
\[
\begin{gathered}
\boldsymbol b^{(0)}
\in\widehat{\mathcal M}_{\mathrm{weak,def}}(L^{(0)}),
\qquad
\boldsymbol b^{(1)}
\in\widehat{\mathcal M}_{\mathrm{weak,def}}(L^{(1)}),\\
\boldsymbol b^{(0)\prime}
\in\widehat{\mathcal M}_{\mathrm{weak,def}}(L^{(0)\prime}),
\qquad
\boldsymbol b^{(1)\prime}
\in\widehat{\mathcal M}_{\mathrm{weak,def}}(L^{(1)\prime}).
\end{gathered}
\]
假设
$\psi_*([\boldsymbol b^{(0)}])
=[\boldsymbol b^{(0)\prime}]$
且
$\psi_*([\boldsymbol b^{(1)}])
=[\boldsymbol b^{(1)\prime}]$。
再假设
\[
\mathfrak{PO}(\boldsymbol b^{(0)})
=\mathfrak{PO}(\boldsymbol b^{(1)}),
\qquad
\pi_{\mathrm{amb}}(\boldsymbol b^{(0)})
=\pi_{\mathrm{amb}}(\boldsymbol b^{(1)}).
\]
那么
\begin{equation}
\mathfrak{PO}(\boldsymbol b^{(0)\prime})
=\mathfrak{PO}(\boldsymbol b^{(1)\prime}),
\qquad
\pi_{\mathrm{amb}}(\boldsymbol b^{(0)\prime})
=\pi_{\mathrm{amb}}(\boldsymbol b^{(1)\prime}),
\tag{5.3.77}\label{eq:5-3-77}
\end{equation}
且存在滤过同构
\[
\begin{aligned}
\psi_*\colon\
&HF\bigl(
(L^{(1)},\boldsymbol b^{(1)}),
(L^{(0)},\boldsymbol b^{(0)});
\Lambda_{0,\mathrm{nov}}
\bigr)\\
&\xrightarrow{\ \cong\ }
HF\bigl(
(L^{(1)\prime},\boldsymbol b^{(1)\prime}),
(L^{(0)\prime},\boldsymbol b^{(0)\prime});
\Lambda_{0,\mathrm{nov}}
\bigr).
\end{aligned}
\]
此外，若
$\{\psi_s\}_{0\le s\le1}$
为辛微分同胚的同痕，
并对每个 $s$ 均满足
$\psi_s(L^{(0)})=L^{(0)\prime}$、
$\psi_s(L^{(1)})=L^{(1)\prime}$，
则 $\psi_{0*}=\psi_{1*}$。
\end{theorem}

\begin{theorem}{4.1.5a}\label{thm:4-1-5a}
设对 $i=0,1$ 有\translatornote{原书实体页 339
在定理 4.1.5a 开头将所有
$\boldsymbol b^{(i)}$
都放在 $L^{(0)}$ 上、
所有 $\boldsymbol b^{(i)\prime}$
都放在 $L^{(1)\prime}$ 上。
这里依后文逐侧的同痕推前条件
分别使用 $L^{(i)}$、
$L^{(i)\prime}$；见疑误 A-0313。}
\[
\boldsymbol b^{(i)}
\in\widehat{\mathcal M}_{\mathrm{weak,def}}(L^{(i)}),
\qquad
\boldsymbol b^{(i)\prime}
\in\widehat{\mathcal M}_{\mathrm{weak,def}}(L^{(i)\prime}),
\]
且
\[
\mathfrak{PO}(\boldsymbol b^{(0)})
=\mathfrak{PO}(\boldsymbol b^{(1)}),
\qquad
\pi_{\mathrm{amb}}(\boldsymbol b^{(0)})
=\pi_{\mathrm{amb}}(\boldsymbol b^{(1)}).
\]
令
$\{\psi_\rho^{(0)}\}_{0\le\rho\le1}$、
$\{\psi_\rho^{(1)}\}_{0\le\rho\le1}$
为哈密顿同痕，满足
\[
\psi_0^{(0)}=\psi_0^{(1)}=\mathrm{id},
\qquad
\psi_1^{(0)}(L^{(0)})=L^{(0)\prime},
\qquad
\psi_1^{(1)}(L^{(1)})=L^{(1)\prime}.
\]
再假设
\[
\psi_{1*}^{(0)}([\boldsymbol b^{(0)}])
=[\boldsymbol b^{(0)\prime}],
\qquad
\psi_{1*}^{(1)}([\boldsymbol b^{(1)}])
=[\boldsymbol b^{(1)\prime}].
\]
那么
\begin{equation}
\mathfrak{PO}(\boldsymbol b^{(0)\prime})
=\mathfrak{PO}(\boldsymbol b^{(1)\prime}),
\qquad
\pi_{\mathrm{amb}}(\boldsymbol b^{(0)\prime})
=\pi_{\mathrm{amb}}(\boldsymbol b^{(1)\prime}).
\tag{5.3.78}\label{eq:5-3-78}
\end{equation}
而且两条同痕诱导弱滤过同构
\[
\begin{aligned}
\bigl(
\{\psi_\rho^{(1)}\}_\rho,
\{\psi_\rho^{(0)}\}_\rho
\bigr)_*\colon\
&HF\bigl(
(L^{(1)},\boldsymbol b^{(1)}),
(L^{(0)},\boldsymbol b^{(0)});
\Lambda_{\mathrm{nov}}
\bigr)\\
&\xrightarrow{\ \cong\ }
HF\bigl(
(L^{(1)\prime},\boldsymbol b^{(1)\prime}),
(L^{(0)\prime},\boldsymbol b^{(0)\prime});
\Lambda_{\mathrm{nov}}
\bigr).
\end{aligned}
\]
此外，若
$\{\psi_{\rho,s}^{(i)}\}_{\rho,s}$
（$i=0,1$）
是哈密顿同痕的双参数族，
且对每个 $s$ 均有
$\psi_{0,s}^{(i)}=\mathrm{id}$、
$\psi_{1,s}^{(i)}(L^{(i)})
=L^{(i)\prime}$，则
\[
\bigl(
\{\psi_{\rho,0}^{(1)}\}_\rho,
\{\psi_{\rho,0}^{(0)}\}_\rho
\bigr)_*
=
\bigl(
\{\psi_{\rho,1}^{(1)}\}_\rho,
\{\psi_{\rho,1}^{(0)}\}_\rho
\bigr)_*.
\]
\end{theorem}

\begin{proof}[定理 4.1.4a 与 4.1.5a 的证明]
(5.3.77) 和 (5.3.78)
由推论 4.6.53 推出。
其余论证与定理 4.1.4、
4.1.5 的证明相同。
为得到“形变”版本，
以定理 5.3.64、5.3.65
替代定理 5.3.1、5.3.14。
“弱”版本则来自定理 5.3.1、
5.3.14、5.3.64、5.3.65
中的同伦等价均为同伦含幺这一事实。
\end{proof}

\begin{proof}[定理 B 的 \textup{(B.1)、(B.3)、(B.4)、(B.5)} 项的证明]
\textup{(B.1)} 是推论 5.2.40 的结果；
\textup{(B.3)}、\textup{(B.4)}
由定理 4.1.4a 推出；
\textup{(B.5)} 由定理 4.1.4a
以及定理 5.3.64 的同伦等价具有函子性推出。
\end{proof}

\begin{proof}[定理 B 的证明收束]
\textup{(B.2)} 已在第 3.8.10 小节证明，
故定理 B 全部得证。
\end{proof}

\begin{proof}[定理 G 的 \textup{(G.1)、(G.3)、(G.4)、(G.5)} 项的证明]
\textup{(G.1)} 由命题 3.7.75 推出；
\textup{(G.3)} 由定理 4.1.4a 推出；
\textup{(G.4)} 由定理 4.1.5a 推出。
\textup{(G.5)} 中双模结构的构造
见第 3.8.10 小节。
从定理 4.1.4a 的同构构造
容易看出，此同构是双模同态。
这就完成了 \textup{(G.5)} 的证明。
\end{proof}

\textup{(G.2)} 将在第 5.4.6 小节证明。
